% Source: https://artofproblemsolving.com/community/c3h2285970_sphere_inscribed_in_right_tetrahedron
% Stable identity: aops-topic-2285970-q01
The largest sphere is the inscribed sphere, tangent to all four faces. Let its centre be $I$ and its radius $r$. Joining $I$ to the four vertices cuts the tetrahedron into four pyramids with apex $I$, one on each face, each of height $r$. Adding their volumes,
\[V=\dfrac{1}{3} rS,\qquad r=\dfrac{3V}{S},\]
where $V$ is the volume and $S$ the total area of the faces. The figure shows the cross-section by the plane $x=y$, through the edge $OZ$ and the midpoint $M$ of the opposite edge; it contains the centre $I$ and is perpendicular to the faces $z=0$ and $x+y+z=1$, so the circle of radius $r$ touches the sides $OM$ and $MZ$.

\begin{asy}
usepackage("amsmath");
size(4cm);
defaultpen(fontsize(10pt));
real r=(3-sqrt(3))/6, h=sqrt(2)/2;
pair O=(0,0), M=(h,0), Z=(0,1), I=(r*sqrt(2),r);
draw(O--M--Z--cycle,linewidth(0.9));
draw(circle(I,r),blue);
dot(I,blue); label("$I$",I,W,blue);
pair T1=(r*sqrt(2),0); pair n=unit((1,h)); pair T2=I+r*n;
draw(I--T1,blue+dashed); draw(I--T2,blue+dashed);
label("$r$",(I+T1)/2,E,blue+fontsize(8pt)); label("$r$",(I+T2)/2,N,blue+fontsize(8pt));
label("$O$",O,SW); label("$M$",M,SE); label("$Z(0,0,1)$",Z,N);
label("$\dfrac{\sqrt2}{2}$",(O+M)/2,S,fontsize(8pt)); label("$1$",(O+Z)/2,W);
\end{asy}

\textbf{Volume.} The base right triangle in the plane $z=0$ has area $\dfrac{1}{2}$ and the height is $1$, so $V=\dfrac{1}{3}\cdot\dfrac{1}{2}\cdot1=\dfrac{1}{6}$.

\textbf{Surface area.} The three faces in the coordinate planes have area $\dfrac{1}{2}$ each. The fourth face is equilateral with side $\sqrt2$ (the distance between any two of $(1,0,0)$, $(0,1,0)$, $(0,0,1)$), so its area is $\dfrac{\sqrt3}{4}\cdot2=\dfrac{\sqrt3}{2}$. Thus $S=\dfrac{3}{2}+\dfrac{\sqrt3}{2}=\dfrac{3+\sqrt3}{2}$.

Hence
\[r=\dfrac{3\cdot\dfrac{1}{6}}{\dfrac{3+\sqrt3}{2}}=\dfrac{1}{3+\sqrt3}=\boxed{\dfrac{3-\sqrt3}{6}}\approx0.211 .\]
Check with coordinates: the centre is $I=(r,r,r)$ (distance $r$ from the three coordinate planes), and its distance to the plane $x+y+z=1$ is $\dfrac{1-3r}{\sqrt3}$, which equals $r$ exactly when $r=\dfrac{1}{3+\sqrt3}$.
